\chapter{Methodology}
\label{chap:methodology}

\section{Research Design}

The project follows an experimental methodology. Its purpose is not only to build attack applications, but to test a defined research question under controlled conditions. The independent change is the attack mechanism or operating condition presented to the IDS. The dependent outcomes are classification accuracy, precision, attack recall, F1, false-positive rate, drift-alert timing, and selected operational network measures. Contiki-NG and Cooja are held constant for the principal attack-family comparison.

The central experiment represents a controlled attack-distribution change. A model is trained on windows from a source attack family and evaluated on windows from a different target attack family. This design tests transfer rather than conventional same-distribution detection. The term concept drift is used carefully. The study induces a change in the relationship between observed features and malicious behaviour by replacing one attack mechanism with another. It does not claim that drift arose naturally in a physical deployment.

\section{Simulation Environment}

The network experiments use Contiki-NG with the Cooja simulator and the RPL-lite routing implementation. The baseline topology contains a DODAG root, normal client nodes, and a designated malicious node. Attack and control configurations use matched seeds so that the comparison is not confounded by unrelated random-state changes. The five seeds are 123456, 234567, 345678, 456789, and 567890.

Each run lasts 540 seconds. An attack application starts in a disabled state, allows the network to establish normal routing behaviour, and activates at 240 seconds. Control applications keep the shared attack flag disabled for the complete run. This delayed activation provides four pre-activation 60-second windows and five post-activation windows in an attack run. It also enables an online monitor to be evaluated against a known change boundary.

\section{Attack Campaign}

The baseline campaign covers nine mechanisms. Blackhole drops forwarded traffic after activation. Grayhole performs selective dropping. Sinkhole advertises the minimum RPL rank while preserving its real internal rank. Increase-rank advertises degraded routing information. DIS flooding repeatedly generates solicitation traffic. DIO suppression interferes with routing advertisement behaviour. Worst-parent deliberately chooses undesirable routing parents. Wormhole introduces a directed radio shortcut between two endpoints. Sybil sends DIO messages using rotating virtual source identities while retaining the underlying mote and RPL state.

The first eight mechanisms reproduce or extend the attack coverage associated with the Gope study. Sybil is the new attack surface in this dissertation because it targets identity perception rather than only forwarding, rank, timing, route choice, or topology. The implementation is described as Sybil-style identity manipulation. It does not model compromise of a cryptographic identity authority.

Five attack and five control runs were validated for every family, giving 90 baseline runs. Validation requires the Cooja completion marker, a radio log, the correct startup behaviour, exactly one delayed activation where relevant, attack effect evidence after activation, and no attack activation in controls. Debug markers establish that code executed, but they are not model features.

\begin{figure}[H]
\centering
\includesvg[width=\textwidth]{figures/attack_coverage}
\caption{Validated Cooja attack coverage. Each family contains five attack and five matched control runs.}
\label{fig:attack-coverage}
\end{figure}

\section{Feature Extraction and Dataset Construction}

Cooja test and radio logs are converted into non-overlapping 60-second windows. A window records provenance, timing, the family and mode, and observed network behaviour. Coarse features describe application transmissions, responses, missed transmissions, reporting nodes, radio transmissions, radio bytes, interference, sender diversity, and selected root or attacker radio activity. Routing features describe DIO, DIS, and DAO events; parent changes; received-DIO rank distributions; low-rank non-root observations; node rank and neighbour state; and inferred topology depth and fan-out.

Two labels are retained for different purposes. The binary run label marks whether a complete run is an attack configuration. The window label is one only when the row belongs to an attack run and its window begins at or after 240 seconds. This distinction allows attack-run provenance to be retained without incorrectly labelling the pre-activation period as malicious.

The frozen baseline dataset contains 810 rows: nine families, ten runs per family, and nine windows per run. The associated manifest contains 90 complete family, mode, and seed groups. An automated audit checks expected families, row counts, window starts, duplicate keys, missing values, numeric parsing, label consistency, and activation flags. The audit passed all checks.

Explicit attack instrumentation is excluded from model inputs. Sixteen fields, including enabled-event counts, drop-event counts, spoofed-DIO counts, and attack-specific effect markers, are classified as validation-only. Metadata such as family, seed, mode, labels, and run ID are also excluded. Fifty-seven generic numeric features remain eligible before a model-specific subset is chosen.

\section{Models and Metrics}

The main experiments use a Gaussian classifier and a CART-style decision tree. These models are intentionally lightweight and transparent. The purpose is to expose transfer and representation failures rather than to optimise a deep model on a fixed benchmark. CART is the principal reported model because it performed more strongly across attack families and its decisions are easier to inspect.

For a binary classifier, precision, recall, and F1 are defined as:
\begin{align}
\mathrm{Precision} &= \frac{TP}{TP+FP}, \\
\mathrm{Recall} &= \frac{TP}{TP+FN}, \\
F1 &= 2\frac{\mathrm{Precision}\times\mathrm{Recall}}{\mathrm{Precision}+\mathrm{Recall}}.
\end{align}
False-positive rate is $FP/(FP+TN)$. Recall and F1 are central because accuracy can remain apparently acceptable when a detector predicts every changed attack window as normal.

\section{Cross-Attack and Adaptation Protocol}

The unit of separation is a complete simulation seed, never an individual CSV row. In a cross-attack test, one family supplies source training evidence and a different family supplies target evaluation evidence. With nine families, CART and Gaussian each produce 72 ordered source-target pairs.

The adaptation experiment augments the source training data with zero, one, two, or three complete target-family seeds. Evaluation uses only target seeds not selected for adaptation. All possible seed combinations are evaluated at each adaptation size where applicable, and mean metrics are reported. This design prevents windows from one simulation appearing in both adaptation and test partitions.

\section{Online Drift Detection}

Two online signals are compared for the focused blackhole-to-sinkhole transition. A one-sided CUSUM monitor consumes generic persistent low-rank RPL state. Its reference distribution is fitted on blackhole windows while excluding the matching seed. A DDM monitor consumes delayed errors produced by the static CART classifier. CUSUM is label-free at decision time, while DDM depends on errors becoming available after labels arrive. Detection rate, delay relative to the 240-second boundary, and control false alarms are recorded.

\section{Robustness and Defence Campaigns}

Robustness is evaluated separately using blackhole, sinkhole, and Sybil under two condition changes. In the attacker-relocated condition, the malicious node is moved from coordinate $(65,75)$ to $(80,75)$. In the lossy-radio condition, UDGM reception success is set to 0.85. Five attack and five control runs are completed for every family-condition combination, giving 60 runs. Static models trained on baseline data are tested on changed conditions, followed by whole-seed target-condition adaptation.

The sinkhole defence campaign contains five controls, five unmitigated attacks, and five defence-enabled runs. The defence identifies a non-root neighbour advertising the root minimum rank and attempts to retain or select a safer parent. Application responses, misses, parent switches, radio transmissions, rank evidence, and defence decisions are measured after activation. Radio transmissions are an activity proxy, not a calibrated energy measure.

\section{External Baseline}

The supplied Gope dataset is analysed separately. It is not merged with the Cooja dataset because the feature schema and run provenance differ. A temporal or source-order split is preferred to random row mixing for the main comparison. Its purpose is to reproduce a strong static baseline and position the Cooja work relative to the motivating paper, not to claim direct external validation of every Cooja result.

\section{Feature-Set Ablation}

Two model-facing feature groups support mechanism analysis. The coarse set includes application transmission, response and miss deltas, reporting-node count, radio transmission count, byte count, interference count, unique senders, selected node radio activity, and basic parent availability. These features represent information that could be collected without detailed RPL rank reconstruction.

The routing set adds DIO, DIS, and DAO counts; parent switches; parent target diversity; significant rank updates; rank distribution summaries; low-rank non-root senders; persistent sender-receiver rank state; node rank and neighbour count; and inferred topology depth and fan-out. Comparing coarse with coarse-plus-routing models tests whether protocol-aware evidence improves a changed mechanism. It is not a pure causal ablation because adding features can also change tree splits and overfitting.

Attack-specific marker columns are never part of either set. For example, the sinkhole advertisement log, blackhole drop marker, Sybil spoof counter, and wormhole endpoint marker are removed. Using these fields could make classification nearly trivial but would answer whether debug code ran rather than whether observable network behaviour indicates attack.

\section{Experimental Questions and Expected Failure Modes}

The static experiment expects in-domain models to outperform cross-attack models. The key null position is that source-family training provides no dependable target-family detection beyond chance or majority behaviour. Evidence against that position would require consistently strong recall and controlled FPR across target mechanisms, not one successful pair.

The adaptation experiment expects performance to improve as complete target seeds are added. A non-monotonic curve is possible because one seed may not represent the target distribution or because new evidence shifts decision boundaries towards false positives. Results are therefore averaged over all available seed combinations at each adaptation size.

The online experiment expects rank-state CUSUM to detect sinkhole earlier than delayed-error DDM because the relevant rank evidence appears immediately after activation while labels arrive later. The robustness experiment expects static thresholds to change under relocation or loss. The defence experiment does not assume success. It tests whether a plausible rank-consistency rule remains safe when connected to routing control.

\section{Validation and Failure Handling}

The campaign was developed iteratively. A run is not counted merely because Cooja started. Validation requires expected logs and behaviour. Early robustness runs without generic RPL logging, builds with damaged compiler flags, shared-source build collisions, weak sinkhole flag binding, and an insufficient timeout were preserved as diagnostics but excluded from final evidence. The final v4 robustness campaign uses isolated code directories, RPL INFO logging, a 541-second timeout, and a strong sinkhole application flag.

Wormhole validation similarly requires tunnel and endpoint radio evidence rather than only an activation message. Sybil validation requires rotating identity events in attack logs and their absence from controls. The defence campaign preserves strict-rejection and alternative-parent diagnostic variants while reporting only the final stable-parent condition as the main prototype.

This failure policy prevents unsuccessful development runs from silently entering the dataset. It also avoids discarding engineering lessons. Diagnostic folders are retained, labelled as non-primary, and discussed where they explain the final design.

\section{Uncertainty and Interpretation}

The campaign has five seeds per attack or condition. Means and complete confusion counts are reported, but five out of five detections do not establish certainty in a wider population. Correlated windows also mean that 45 windows per mode are not 45 independent networks. The run is the primary experimental unit even when the classifier operates on windows.

The adaptation summaries average over seed combinations, which reduces dependence on one convenient split. The cross-attack matrix reports every ordered family pair. Robustness results report FPR alongside F1 so that an always-alerting model is not mistaken for a robust detector. Negative findings are retained when metrics worsen.

\section{Ethical and Reproducibility Considerations}

All attacks run in an isolated simulator against project-owned configurations. No public network, third-party device, credential, or personal dataset is targeted. The work is defensive security research intended to understand IDS limitations. Attack code and reproducibility artefacts should be handled in accordance with university guidance and described with enough detail for academic evaluation without claiming deployment authorisation.

Reproducibility is supported through fixed seeds, delayed activation, preserved configurations, raw logs, generated CSVs, manifests, audit reports, source tables for figures, and local scripts. No external API is required for the core experiments. The exact Contiki-NG revision and local environment should be recorded in the final submission metadata before release.

\section{Analysis Workflow}

The analysis proceeds in a fixed order. First, family validators establish that every selected run completed and that attack/control behaviour matches its configuration. Second, the common extractor creates windows and the frozen-data audit verifies integrity. Third, static in-domain and cross-attack models are evaluated without markers. Fourth, target-seed combinations are added for adaptation. Fifth, chronological sinkhole windows are evaluated by CUSUM and DDM. Sixth, baseline models are applied to the separate robustness conditions and adapted with condition-specific seeds. Finally, the defence campaign is summarised independently.

This ordering prevents later experiments from redefining earlier evidence. The baseline dataset remains frozen when robustness or defence runs are added. Figure generation reads result CSVs rather than rerunning models. The dissertation can therefore trace every claim to a stable source table.

The core argument is evaluated through triangulation. Static cross-attack matrices establish failure. Adaptation curves establish recoverability. Online monitors establish an observable change point for one focused transition. Robustness runs challenge environmental stability. The defence result tests the boundary between detection and control. Agreement is not required across every experiment; disagreement identifies where a feature or method is mechanism-specific.
